\documentclass{article}
\usepackage[T1]{fontenc}
\usepackage{lmodern}
\usepackage{graphicx}
\usepackage{amsmath,amssymb}
\usepackage[a4paper,margin=2cm]{geometry}
\usepackage[absolute,overlay]{textpos}
\setlength{\TPHorizModule}{1mm}
\setlength{\TPVertModule}{1mm}
\usepackage[indonesian]{babel}
\usepackage{hyperref}
\setlength{\emergencystretch}{2em}
% unit: Halaman 4

\begin{document}

% === Halaman 4 (python/native) ===

• E dan D adalah fungsi enkripsi dan dekripsi dari suatu cipher blok dengan 

parameter kunci K:



C = EK(P) 

dan P = DK(C)

• Dimisalkan contoh sebuah fungsi E yang sederhana, tapi tidak cukup aman, sbb:

• D adalah kebalikan (invers) dari E. Contoh fungsi D yang sederhana adalah sbb:



4

1.  XOR-kan blok plainteks P dengan K 

   2.  lalu geser secara wrapping bit-bit dari hasil langkah 1 satu posisi ke kiri


C = EK(P) = (P  K) << 1

1. Geser secara wrapping  bit-bit ciphertes C  satu posisi ke kanan

   2. Lalu XOR-kan blok hasil Langkah 1 dengan  K 


P = DK(C) = (C  >> 1)  K

\end{document}
